\documentclass[11pt,a4paper]{article}
\usepackage[margin=0.85in,top=1in,bottom=1in]{geometry}
\usepackage{amsmath,amssymb,amsfonts}
\usepackage{graphicx}
\usepackage{float}
\usepackage{enumitem}
\usepackage{microtype}
\usepackage{caption}
\usepackage[hidelinks]{hyperref}
\usepackage[most,breakable]{tcolorbox}
\tcbuselibrary{breakable,skins}
\usepackage[utf8]{inputenc}
\usepackage[T1]{fontenc}
\usepackage{lmodern}
\captionsetup{font=small,labelfont=bf,skip=4pt}
\setlist{leftmargin=1.5em,topsep=2pt}

%% Breakable card — prevents content cutoff across pages
\newtcolorbox{solcard}[3]{
  breakable, enhanced,
  colback=white, colframe=gray!50, arc=2pt,
  title={\textbf{#1}\hfill\textsf{\small #3\,pts}},
  fonttitle=\normalsize\sffamily\bfseries,
  before upper={\small\textit{#2}\par\smallskip\hrule height 0.4pt\smallskip},
  top=4pt, bottom=6pt, left=7pt, right=7pt,
  parbox=false,
}

\title{\textbf{Josephson effect --- Solution}}
\author{\normalsize X25 (2025) — English translation}
\date{}
\begin{document}\maketitle\vspace{-0.5em}
\begin{solcard}{A1}{Let at the initial moment of time $t = 0$ the current through the contact is equal to to zero. At this moment, voltage $V$ was applied to the contact.  Determine the dependence of the current on time $I(t)$. Express your answer using $\Phi_0$, $V$, $I_c$, $R$.}{0.50}

The voltage sets the derivative of the phase difference at the contact $$ \dot{\varphi} = \frac{2 \pi} {\Phi_0}V. $$Since before the voltage was turned on, the current equals zero, initial phase $\varphi(0) = 0$. (For switching on voltage phase not can change instantaneously, since fast change in phase leads to infinitely large voltage.) Then phase have form $$ \varphi = \frac{2 \pi} {\Phi_0}V t. $$Then total current consists of the normal current $I_n = V/R$ and the superconducting current $I_s = I_c \sin\varphi$

\par\smallskip\noindent\textbf{Answer:} $$I = \dfrac{V}{R} + I_c \sin \left( \dfrac{2\pi}{\Phi_0}V t\right)$$
\end{solcard}

\begin{solcard}{A2}{Let a direct current $I > I_c$ be passed through the contact.  Write down the differential equation for $\varphi(t)$.}{0.50}

Let us write the total current as the sum of the normal current and the superconducting current, expressing the voltage through the derivative of the phase, as a result we obtain

\par\smallskip\noindent\textbf{Answer:} $$I = I_c \sin \varphi + \dfrac{\Phi_0}{2 \pi R} \dot{\varphi}$$
\end{solcard}

\begin{solcard}{A3}{Differentiate the equation with respect to time. Eliminate the phases from the equation and obtain an expression for $\dot{V}(t)^2$ in terms of $V$, $I$ and the system parameters.}{0.70}

Differentiating the equation, we obtain $$ (I_c \cos \varphi )\dot{\varphi} + \frac{\Phi_0}{2 \pi R} \ddot{\varphi} = 0. $$Derivative $\dot{\varphi}$ express through voltage $$ V (I_c \cos \varphi) + \frac{\Phi_0}{2 \pi R} \dot{V} = 0, $$and the cosine of the phase from the sine using the original equation, $$ I_c \sin \varphi = I - \frac{V}{R}, \quad I_c^2 \cos^2 \varphi = I_c^2 - \left(I - \frac{V}{R} \right)^2, $$Finally we obtain

\par\smallskip\noindent\textbf{Answer:} $$ \dot{V}^2 = \left(\frac{2 \pi R}{\Phi_0} \right)^2 V^2 \left(I_c^2 -  \left(I - \frac{V}{R} \right)^2\right) $$
\end{solcard}

\begin{solcard}{A4}{In the previous equation, go to the variable $u = 1/V$, express $\dot{u}^2$ in terms of $u$, $I$ and the system parameters.}{0.30}

Let us express the derivative of the voltage $$ \dot{V} = - \frac{\dot{u}}{u^2}, $$substitute in equation $$ \frac{\dot{u}^2}{u^4} = \left(\frac{2 \pi R}{\Phi_0} \right)^2 \frac{1}{u^2} \left(I_c^2 -  \left(I - \frac{1}{R u} \right)^2\right). $$Multiplying by $u^4$, we obtain

\par\smallskip\noindent\textbf{Answer:} $$ \dot{u}^2 = \left(\frac{2 \pi R}{\Phi_0} \right)^2\left(I_c^2 u^2 - \left(I u - \frac{1}{R} \right) ^2\right) $$
\end{solcard}

\begin{solcard}{A5}{You should get an equation similar to the law of conservation of energy for oscillations of a harmonic oscillator. Determine the frequency of oscillations $\omega$, the value $u_0$ about which the oscillations occur, and the amplitude of oscillations $A$.}{0.80}

Differentiating the equation from the previous paragraph, we obtain the vibration equation $$ \ddot{u} = -  \left(\frac{2 \pi R}{\Phi_0} \right)^2 \left( (I^2 - I_c^2) u  - \frac{I}{R} \right). $$From coefficient before $u$ find frequency oscillation $$ \omega = \frac{2\pi R}{\Phi_0} \sqrt{I^2 - I_0^2}, $$and also position equilibrium $$ u_0 = \frac{I}{R (I^2 - I_c^2)}. $$For that, so that find amplitude oscillation, determinable value $u$, for which $\dot{u} = 0$, obtain square equation $$ (I^2 -I_c^2)u^2 - \frac{2I}{R} u + \frac{1}{R^2} = 0, $$root which $$ u = \frac{1}{I^2 - I_c^2} \left( \frac{I}{R} \pm \frac{I_c}{R}\right), \quad u_1 = \frac{1}{R(I + I_c)}, \quad u_2 = \frac{1}{R(I-I_c)}. $$That same result one can find, determine maximum and minimum value voltage from the equation $$ I = I_c \sin \varphi + \frac{V}{R}, \quad V_{min} = R(I- I_c), \quad V_{max} = R(I + I_c). $$Then amplitude $$A = \frac{1}{2} (u_2 - u_1) = \frac{I_c}{R(I^2 - I_c^2)}.$$

\par\smallskip\noindent\textbf{Answer:} $$ \omega = \frac{2\pi R}{\Phi_0} \sqrt{I^2 - I_0^2}; \quad u_0 = \frac{I}{R (I^2 - I_c^2)}; \quad A =  \frac{I_c}{R(I^2 - I_c^2)}. $$
\end{solcard}

\begin{solcard}{A6}{Write down the dependence $V(t)$. Consider that at $t = 0$ the value $V(t)$ has a minimum.}{0.20}

The dependence of the quantity $u$ on time $$ u = u_0 + A \cos \omega t, $$where phase be determined the more condition, that $V$ minimum in initial moment time, and means $u$ is maximum.

\par\smallskip\noindent\textbf{Answer:} $$ V = \frac{R(I^2 - I_c^2)}{I + I_c \cos \left(\dfrac{2\pi R}{\Phi_0} \sqrt{I^2 - I_0^2} t \right)} $$
\end{solcard}

\begin{solcard}{A7}{Construct high-quality graphs of the dependence $V(t)$ for cases $I = 1.1I_c$ and $I= 3 I_c$, indicate the coordinates of the characteristic points.}{0.60}
\par\smallskip\noindent\textbf{Answer:} Graph for case $I = 1.1 I_c$
\begin{figure}[H]
  \centering
  \includegraphics[width=0.72\textwidth]{T3_S_EN_assets/fig1.jpg}
\end{figure}
\begin{figure}[H]
  \centering
  \includegraphics[width=0.72\textwidth]{T3_S_EN_assets/fig2.jpg}
\end{figure}
\end{solcard}

\begin{solcard}{A8}{Find the average (over time) voltage across the contact. Express your answer in terms of $R$, $I$, $I_c$.}{0.80}

Let's express the voltage through the phase derivative and integrate over time $$ V =\frac{\Phi_0}{2\pi} \dot{\varphi}, \quad \int\limits_{0}^{T} V(t) dt  = \frac{\Phi_0}{2\pi} (\varphi(T) - \varphi(0)). $$If $T$ – period oscillation, that $\varphi(T) - \varphi(0) = 2\pi$, and mean value voltage equal $$ \langle V \rangle =  \frac{1}{T} \int\limits_{0}^{T} V(t) dt = \frac{\Phi_0}{T} = \frac{\Phi_0}{2\pi} \omega = R\sqrt{I^2 -I_0^2}. $$This the same result can be obtained by integrating the explicit expression for $V(t)$ $$ \langle V \rangle = \frac{1}{T} \int\limits_{0}^{T} V(t) dt = \frac{1}{T} \int\limits_0^{T} \frac{R(I^2 - I_c^2)}{I + I_c \cos \left(\omega t \right)} dt = \frac{1}{2\pi} \int\limits_0^{2\pi} \frac{R(I^2 - I_c^2)}{I + I_c \cos \theta} d\theta, $$where $\theta = \omega t$. This integral one can calculate  with means substitution $$ \xi = \tan \frac{\theta}{2} , \; \cos \theta = \frac{1-t^2}{1+ t^2}, \; d\theta = \frac{2 dt}{1+ t^2}. $$$$ \int\limits_0^{2\pi} \frac{d\theta}{I+I_c \cos \theta} = \int\limits_{-\pi}^{\pi} \frac{d\theta}{I+I_c \cos \theta} = \int\limits_{-\infty}^{\infty} \frac{\frac{2 dt}{1+t^2}}{I + I_c \frac{1-t^2}{1+t^2}} = 2 \int\limits_{-\infty}^{\infty} \frac{dt}{I+ I_c + (I-I_c)t^2} = \frac{2\pi}{\sqrt{I^2 - I_c^2}}. $$This gives the same answer

\par\smallskip\noindent\textbf{Answer:} $$ V = R \sqrt{I^2 - I_c^2} $$
\end{solcard}

\begin{solcard}{A9}{Determine the dependence of the energy $E(\varphi)$ of the Josephson junction as a function of the phase difference $\varphi$ on it. The answer may also include $\Phi_0$ and $I_c$. Consider that $E(0) = 0$.}{0.70}

Source power $$P = IV = I_s V = I_c \sin\varphi \dfrac{\Phi_0 }{2\pi} \dot{\varphi}.$$The contribution of the normal current to the power has the form $$ P_n = I_n V = \frac{1}{R} \left( \frac{\Phi_0}{2\pi}\right)^2 \dot{\varphi}^2. $$For magnification characteristic time change phase $T$ this power vary as $1/T^2$, therefore the total released heat behaves as $(1/T^2)T \sim 1/T$ and it can be made arbitrarily small by increasing $T$. Integrating in time, we find the work done by the the source. $$ A = \frac{\Phi_0 I_c}{2\pi} \int \sin \varphi \dot{\varphi} dt =  \frac{\Phi_0 I_c}{2\pi} \int\limits_{0}^{\varphi} \sin \varphi d \varphi =  \frac{\Phi_0 I_c}{2\pi} (1 - \cos \varphi). $$This work is equal to the energy of the Josephson junction.

\par\smallskip\noindent\textbf{Answer:} $$ E =  \frac{\Phi_0 I_c}{2\pi} (1 - \cos \varphi) $$
\end{solcard}

\begin{solcard}{A10}{For a typical Josephson junction, the value of the critical current is $I_c = 0.5~\mathrm{mA}$, the resistance is $R = 0.7 ~\mathrm{\Omega}$. Fundamental constants $e = 1.602 \times 10^{-19}~\mathrm{C}$, $\hbar = 1.055 \times 10^{-34}~\mathrm{J} \cdot \mathrm{s}$. Find the numerical value of the energy of the Josephson junction $E(\pi)$ and the frequency of voltage oscillations at constant current through the contact $I = 2 I_c$.}{0.20}

$$ E(\pi) = \frac{\hbar I_c}{e} = 3.3 \times 10^{-19}~\mathrm{J}, $$$$ \omega = 2 \sqrt{3} \frac{e R I_c}{\hbar} = 1.84 \times 10^{12}~\mathrm{s}^{-1} $$

\end{solcard}

\begin{solcard}{B1}{Let a direct current $I$ flow into the contact (see figure in the introduction to the problem). Write down a differential equation describing the change in phase difference $\varphi$. The answer may also include $\Phi_0$, $R$, $C$, $I_c$.}{0.30}

To the expression for the current from part A2, you need to add the current that goes to charge the capacitor, $I_C = \dot{q} = C \dot{V}$, and the voltage derivative is expressed through the second derivative of the phase

\par\smallskip\noindent\textbf{Answer:} $$ I = I_c \sin \varphi + \frac{\Phi_0}{2\pi R} \dot{\varphi} + \frac{\Phi_0 C}{2\pi} \ddot{\varphi} $$
\end{solcard}

\begin{solcard}{B2}{Let the external current be zero and the phase difference $\varphi \ll 1$. In this case, $\varphi$ (and therefore other quantities characterizing the contact) will perform damped oscillations. Determine the vibration frequency $\omega_p$.}{0.60}

For small phase values, the equation has the form (after dividing by the coefficient before the second derivative) $$ \ddot{\varphi} + \frac{1}{RC} \dot{\varphi} + \frac{2 \pi I_c }{\Phi_0 C}\varphi = 0. $$This equation have standard form equation of damped oscillation $$ \ddot{\varphi} + 2 \gamma \dot{\varphi}  + \omega_0^2 \varphi = 0, $$with coefficient $$ \omega_0^2 =  \frac{2 \pi I_c }{\Phi_0 C}, \quad \gamma = \frac{1}{2 RC}. $$Then frequency $\omega = \sqrt{\omega_0^2 - \gamma^2}$, This answer one can obtain, substitute in equation solution form $\varphi = A e^{i \omega t}$, find corresponding value $\omega$ and taking the real part.

\par\smallskip\noindent\textbf{Answer:} $$ \omega_p = \sqrt{ \frac{2 \pi I_c }{\Phi_0 C}- \frac{1}{4 R^2 C^2}} $$
\end{solcard}

\begin{solcard}{B3}{Determine the quality factor of the oscillatory system from the previous paragraph $Q$. At what values of capacitance in the system are oscillations possible rather than damped motion?}{0.60}

Quality factor can be expressed as $$ Q = \frac{\omega_0}{2 \gamma} = R \sqrt{\frac{2 eI_c C}{\hbar}}. $$Oscillation possible, while the value of $\omega_p$ is real; otherwise the solution of the equation for $\varphi$ be represent itself sum of damped exponential, hence find $$ C \ge \frac{\Phi_0}{8 \pi R^2 I_c}. $$

\par\smallskip\noindent\textbf{Answer:} $$ Q = R \sqrt{\frac{2 eI_c C}{\hbar}}, \quad C \ge \frac{\Phi_0}{8 \pi R^2 I_c}. $$
\end{solcard}

\begin{solcard}{C1}{Let the critical currents through the contacts be $I_{ac}$ and $I_{bc}$ respectively. Then the maximum superconducting current $I$ that can flow through such a contact will depend on the magnetic flux through the interferometer. Determine this maximum value $I_0$, as well as the cosines of the phase differences $\varphi_a$, $\varphi_b$ at which this maximum is achieved. Express your answer using $I_{ac}$, $I_{bc}$, $\Phi$, $\Phi_0$. Your answers should not contain inverse trigonometric functions.}{1.20}

The total current through the contacts is $$ I = I_a + I_b = I_{ac} \sin \varphi_{a} + I_{bc} \sin \varphi_{b}. $$Use relation between phase $\varphi_b = \varphi_a - \Delta \varphi$, where $\Delta \varphi = 2\pi \Phi/\Phi_0$, obtain after transformation $$ I = I_{ac} \sin \varphi_a + I_{bc} \sin (\varphi_a - \Delta \varphi) = (I_{ac} + I_{bc} \cos \Delta \varphi) \sin \varphi_a  - I_{bc} \sin \Delta \varphi \cos \varphi_a = I_0 \sin (\varphi_a - \psi), $$where amplitude $$ I_0^2 = (I_{ac} + I_{bc} \cos \Delta \varphi)^2 + (I_{bc} \sin\Delta \varphi)^2 = I_{ac}^2 + I_{bc}^2 + 2 I_{ac} I_{bc} \cos \Delta \varphi, $$and phase be determined relation $$ I_0 \cos \psi = I_{ac} + I_{bc} \cos \Delta \varphi, \quad I_0 \sin \psi = I_{bc} \sin \Delta \varphi. $$Current be maximum, if argument sine equal $\pi/2$, that is $\varphi_a = \psi + \pi/2$, from which $$ \cos \varphi_a  = - \sin \psi= - \frac{I_{bc} \sin \Delta \varphi}{\sqrt{ I_{ac}^2 + I_{bc}^2 + 2 I_{ac} I_{bc} \cos \Delta \varphi}}. $$The cosine of the second phase can be obtained using the symmetry between the contacts $a$ and $b$ (if we write expression for $I_0$ through $\varphi_b$, the expression differs only by the substitution of $a$ for $b$ in the indices and change sign $\Delta \varphi$). Otherwise it can be found by direct calculation $$ \cos \varphi_b = \cos (\varphi_a - \Delta \varphi) = \cos \varphi_a \cos \Delta \varphi + \sin \varphi_a \sin \Delta \varphi. $$Substitute expression for cosine and sine through $\psi$, obtain $$ \cos \varphi_b = - \sin \psi \cos \Delta \varphi+ \cos \psi \sin \Delta \varphi = - \frac{I_{bc} \sin \Delta \varphi}{I_0} \cos \Delta \varphi + \frac{I_{ac} + I_{bc} \cos \Delta \varphi}{I_0} \sin \Delta \varphi, $$$$ \cos \varphi_b = \frac{I_{ac} \sin \Delta \varphi}{\sqrt{ I_{ac}^2 + I_{bc}^2 + 2 I_{ac} I_{bc} \cos \Delta \varphi}}. $$

\par\smallskip\noindent\textbf{Answer:} $$ I_0 = \sqrt{I_{ac}^2 + I_{bc}^2 + 2 I_{ac} I_{bc} \cos \left( 2\pi \frac{\Phi}{\Phi_0}\right)} $$$$ \cos \varphi_a = - \frac{I_{bc} \sin \left( 2\pi \frac{\Phi}{\Phi_0}\right)}{ \sqrt{I_{ac}^2 + I_{bc}^2 + 2 I_{ac} I_{bc} \cos \left( 2\pi \frac{\Phi}{\Phi_0}\right)}} $$$$ \cos \varphi_b =  \frac{I_{ac} \sin \left( 2\pi \frac{\Phi}{\Phi_0}\right)}{ \sqrt{I_{ac}^2 + I_{bc}^2 + 2 I_{ac} I_{bc} \cos \left( 2\pi \frac{\Phi}{\Phi_0}\right)}} $$
\end{solcard}

\end{document}